\documentclass[11pt,a4paper]{article}
\usepackage[T1]{fontenc}
\usepackage{lmodern,amsmath,amssymb,amsthm,mathtools,booktabs,microtype}
\usepackage[margin=27mm]{geometry}
\usepackage[colorlinks=true,linkcolor=blue,citecolor=blue,urlcolor=blue]{hyperref}
\setlength{\emergencystretch}{2em}
\newtheorem{theorem}{Theorem}[section]
\newtheorem{lemma}[theorem]{Lemma}
\newtheorem{corollary}[theorem]{Corollary}
\newcommand{\Q}{\mathbb Q}\newcommand{\Z}{\mathbb Z}
\DeclareMathOperator{\Tr}{Tr}
\DeclareMathOperator{\N}{N}
\title{Polynomial complement levels in local elliptic detectors}
\author{}\date{16 November 2026}
\begin{document}\maketitle
\begin{abstract}
We prove a local extension of the fixed-level complement criterion.
A polynomial level of degree $2m$ may have negative valuation at an
input denominator. Weighted comparison and reciprocal-abscissa tests
compensate for that loss. The comparison exponent may be any integer
$s$ with $2s\ge m$, and the reference exponent any even integer
$r\ge2+m/2$. At an odd positive contact of the physical weight,
the resulting equation excludes input denominators and gives the
required elliptic index congruence. The result concerns the guarded
point locus and is independent of a global realization theorem.
\end{abstract}

\section{Local data}

Let $K=\Q(\sqrt2)$, with nonidentity automorphism $\sigma$, and let
\[
             E:\ y^2=x^3-Dx,\qquad D=(3-\sqrt2)^2.
\]
Fix an odd prime $p>3$ of good reduction at both embeddings, with
$D$ a unit. Normalize valuations over $p$ by $v(p)=1$; these valuations
are integral since $K/\Q$ is unramified at $p$.
Let $P,H,U,A$ be $K$-points, with $P$ non-torsion. Write
\[
                  H_0=H-P,\qquad V=U-A.
\]
Assume all points occurring in the following formulas have nonzero
abscissa and ordinate and all difference denominators are nonzero.
These are the ordinary guarded point-chart conditions.
Put
\[
 f(Q)=1/x(Q),\quad z(Q)=-x(Q)/y(Q),\quad
 L=z(U)/z(H),\quad L_0=z(V)/z(H_0).
\]

Assume the source supplies indices $h,u,n\in\frac12\Z$ satisfying
\begin{equation}\label{eq:indices}
 [2]H=[2h]P,\qquad [2]U=[2u]P,\qquad [2]A=[2n]P.
\end{equation}
At odd $p$ these indices are $p$-integral. The factor two allows both
the integral and half-integral label conventions. The non-torsion
source charts exclude zero indices for the reference points.

Take $\beta=X$ or $X/2$, with $X\in\Q^\times$, and a nonzero polynomial
level
\begin{equation}\label{eq:level}
 a(X)=cM(X)^2,\qquad M\in\Z_{(p)}[X],\qquad \deg M=m,\qquad
 v_p(c)\in2\Z_{\ge0}.
\end{equation}
Here $\Z_{(p)}$ denotes the rational numbers integral at $p$.
Choose integers $s\ge0$ and $r\ge0$ with
\begin{equation}\label{eq:exponents}
                  2s\ge m,\qquad r\text{ even},\qquad r\ge2+m/2.
\end{equation}
Define
\[
 C=\beta^s\left(\frac L\beta-\frac{\sigma L_0}{\beta}\right),
 \qquad I=\beta^s\left(\frac L\beta-1\right),
\]
\[
 J=\beta^r f(H)f(U),\qquad
 J'=\beta^r\sigma\bigl(f(H_0)f(V)\bigr).
\]
Let $w\in\Q^\times$ have positive odd valuation
$\nu=v_p(w)$. Set
\begin{equation}\label{eq:detector}
 D_1=1+\frac{C^2}{w}+\frac{J^2}{w^2},\qquad
 D_2=1+\frac{I^2}{w}+\frac{(J')^2}{w^2},\qquad W=D_1D_2.
\end{equation}
Require $\N(W)\Tr(W)\ne0$, and put
\[
                  \tau=\frac{\N(W)}{\Tr(W)},\qquad b=w\tau^2.
\]
Multiplying $b$ or $w$ by a fixed $p$-adic unit does not affect the
valuation conclusions below.

\begin{theorem}\label{thm:local}
Under the preceding hypotheses, if
\begin{equation}\label{eq:complement}
               q^2=d\,a(X)(a(X)-b),\qquad v_p(d)\in2\Z,
\end{equation}
has a rational solution, then
\[
                         v_p(\beta)\ge0,\qquad
                         \beta(n-\beta)\in p\Z_{(p)}.
\]
In particular, if $v_p(\beta)=0$, then $n\equiv\beta\pmod p$.
\end{theorem}

\section{Valuations of the point tests}

\begin{lemma}\label{lem:table}
For a guarded non-torsion point $Q$, the following table gives all
possible valuation patterns, with $e>0$.
\[
\begin{array}{c|cc}
\text{reduction of }Q&v(z(Q))&v(f(Q))\\ \hline
O&e&2e\\
(0,0)&e&-2e\\
\text{nonzero-abscissa two-torsion}&-e&0\\
\text{other finite point}&0&0
\end{array}
\]
If $v(z(U)/z(H))=-k$ with $k\ge0$, then on every branch other
than $O/O$,
\begin{equation}\label{eq:wrong}
                         v(f(H)f(U))\le2k.
\end{equation}
\end{lemma}
\begin{proof}
At $O$, the usual formal parameter gives $x=z^{-2}$ times a unit.
Near $(0,0)$ the equation and $D$ being a unit give
$v(x)=2v(y)>0$, hence $v(z)=v(y)$ and $v(f)=-2v(y)$.
Near another two-torsion point, $x$ is a unit and $y$ vanishes.
At the remaining finite points both are units.

Write $a=v(z(H))$ and $v(z(U))=a-k$. If $H$ reduces to $O$
and $U$ reduces to $(0,0)$, the reciprocal product has valuation
$2a-2(a-k)=2k$. If $U$ is ordinary it has valuation $2k$;
if $U$ approaches another two-torsion point, then $0<a<k$ and
the product has valuation $2a<2k$.
If $H$ approaches $(0,0)$ and $U$ approaches $O$, the product
has valuation $-2k$; all other cases with this reference have
no larger value. An ordinary reference or a reference approaching
another two-torsion point gives product valuation at most zero
when the numerator parameter has valuation $a-k$.
This proves~\eqref{eq:wrong}, including $k=0$.
\end{proof}

In particular, reciprocal abscissas have even valuations. Because $r$
is even, $v(J)$ and $v(J')$ are even. The three finite summand
valuations in either factor~\eqref{eq:detector} have respective classes
\[
                       0\pmod4,\qquad 1\pmod2,\qquad2\pmod4.
\]
Their minimum is unique; zero summands may simply be omitted.
Thus each factor has valuation at most zero at each place over $p$.

\begin{lemma}\label{lem:trace}
For some embedding over $p$,
\[
                              v(W)\ge v_p(\tau).
\]
At such an embedding, a strict lower bound on $v_p(\tau)$
is a strict lower bound on each factor and each summand
in~\eqref{eq:detector}.
\end{lemma}
\begin{proof}
Since $\tau=1/\Tr(W^{-1})$, the ultrametric inequality gives
\[
 v_p(\tau)\le
       \max\{v(W),v(\sigma W)\}.
\]
Both factor valuations are nonpositive, so each is at least their
sum $v(W)$. The unique-minimum property then gives the same lower
bound for every summand.
\end{proof}

\section{Proof of the local theorem}

Suppose first that $v_p(\beta)=-k<0$. Since $p$ is odd,
$v_p(X)=-k$. Equation~\eqref{eq:level} gives
\[
                     v_p(a(X))\ge-2mk,
\]
and this valuation is even. On the other hand,
$v_p(b)=\nu+2v_p(\tau)$ is odd. If it were smaller than
$v_p(a(X))$, the right side of~\eqref{eq:complement} would
have odd valuation. Equality is excluded by parity. Hence
\begin{equation}\label{eq:tau}
                         v_p(\tau)>-mk-\nu/2.
\end{equation}
At the embedding of Lemma~\ref{lem:trace}, the comparison summands give
\[
 v(C),v(I)>-mk/2+\nu/4.
\]
After dividing by $\beta^s$, the two normalized comparison errors
have valuations greater than
\[
                           (s-m/2)k+\nu/4>0.
\]
Consequently $L/\beta$ and $\sigma L_0/\beta$ are units congruent
to one. Both raw ratios have valuation $-k$.
The reference summands give
\begin{equation}\label{eq:J}
                        v(J),v(J')>-mk/2+3\nu/4.
\end{equation}

On any wrong branch of either reference pair,
Lemma~\ref{lem:table} gives
\[
                 v(\beta^r f(H)f(U))\le(2-r)k\le-mk/2,
\]
contradicting~\eqref{eq:J}. Thus $H,U$ are formal at $O$ at one
embedding, and $H_0,V$ are formal at $O$ at the conjugate embedding.

The formal logarithm is an isometry. Equations~\eqref{eq:indices}
therefore give the valuation identities
\[
 v\!\left(\frac{z(U)}{z(H)}\right)=v_p(u/h),\qquad
 v\!\left(\frac{z(V)}{z(H_0)}\right)
                              =v_p((u-n)/(h-1)).
\]
One can obtain these identities by cross-multiplying the indices
and applying the formal logarithm; multiplication by two is
invertible in the formal group. Both ratios have negative valuation.
Since $u,u-n$ are $p$-integral, this forces both $h$ and $h-1$
to have positive valuation, which is impossible. This proves
$v_p(\beta)\ge0$.

If $v_p(\beta)=0$, the level is integral with even valuation, and
the complement excludes every negative odd valuation of $b$.
The same argument with $k=0$ gives positive normalized comparison
errors and positive reference-product valuations. The table again
forces both reference pairs to be formal at the corresponding
embeddings. On the first formal neighborhood, the formal logarithm
agrees with $z$ up to a multiplicative factor in $1+p\mathcal O$.
Multiplying the resulting congruences by the integral indices gives
\[
                    u\equiv\beta h,\qquad
                    u-n\equiv\beta(h-1)\pmod p.
\]
Subtracting proves $n\equiv\beta\pmod p$.
Finally, if $v_p(\beta)>0$, both $n$ and $\beta$ are integral at
$p$, so $\beta(n-\beta)$ is already divisible by $p$.
This completes the proof of Theorem~\ref{thm:local}.

\begin{corollary}
For $a(X)=17(X^2+B)^2$, with $B$ a positive integer, one may take
$s=1$ and $r=4$. The tests are then
\[
 C=L-\sigma L_0,\quad I=L-\beta,\quad
 J=\beta^4f(H)f(U),\quad
 J'=\beta^4\sigma(f(H_0)f(V)).
\]
\end{corollary}

\section{Scope}

The theorem applies on the nonzero-ordinate guarded source locus.
It makes no assertion about extending a quadratic chart through an
exact torsion point. It also does not provide a global height bound,
real preparation, a rational realization theorem, or a new defining
variable count. These requirements remain separate from the proved
local implication.

\begin{thebibliography}{9}
\bibitem{FiniteClasses}
\emph{Finite square-class conditions for native detector inputs},
project source, 16 November 2026,
\nolinkurl{papers/pole_parameters/finite_complement_classes_20261116_v46.tex}.
\bibitem{Quartic}
\emph{A quartic in 74 variables from weighted point products},
project source, 16 November 2026,
\nolinkurl{papers/quartic_existence/quartic_abc74_20261116_v46.tex}.
\bibitem{Direct}
\emph{Thirteen Variables from Direct Elliptic Index Comparison},
project source, 10 November 2026,
\nolinkurl{papers/conditional_few_variables/direct_input13_20261110_v22.tex}.
\end{thebibliography}
\end{document}
